\documentclass[11pt,a4paper]{article}
\usepackage[margin=25mm,headheight=14pt]{geometry}
\usepackage{fontspec}
\setmainfont{lmroman10-regular.otf}[BoldFont=lmroman10-bold.otf,ItalicFont=lmroman10-italic.otf,BoldItalicFont=lmroman10-bolditalic.otf]
\setsansfont{lmsans10-regular.otf}
\setmonofont{lmmono10-regular.otf}[Scale=MatchLowercase]
\newfontfamily\authorsymbol{Noto Sans Symbols}[Scale=MatchUppercase]
\usepackage{amsmath,amssymb,amsthm,mathtools,booktabs}
\usepackage{microtype}
\usepackage[hidelinks,unicode]{hyperref}
\usepackage{fancyhdr}
\pagestyle{fancy}\fancyhf{}
\fancyhead[L]{\small\itshape Ordered carry frontiers for parallel addition}
\fancyhead[R]{\small\thepage}\renewcommand{\headrulewidth}{0pt}
\setlength{\parindent}{1.2em}\setlength{\parskip}{0pt}
\setlength{\emergencystretch}{2em}
\newtheorem{theorem}{Theorem}
\newtheorem{proposition}[theorem]{Proposition}
\newtheorem{lemma}[theorem]{Lemma}
\newtheorem{corollary}[theorem]{Corollary}
\newcommand{\ind}{\mathbf 1}
\newcommand{\bits}{\{0,1\}}
\newcommand{\andb}{\mathbin{\&}}
\DeclareMathOperator{\rank}{rank}
\hypersetup{pdftitle={Ordered Carry Frontiers for Parallel Addition},pdfauthor={lmtlssss🜐},pdfsubject={Mathematics: arithmetic automata and exact counting},pdfkeywords={shared input, carry states, finite automata, exact counting, masked addition}}
\title{\LARGE Ordered Carry Frontiers\\for Parallel Addition}
\author{lmtlssss{\authorsymbol 🜐}}
\date{27 October 2026\\[3pt]\small\href{https://doi.org/10.5281/zenodo.23702847}{10.5281/zenodo.23702847}}
\begin{document}
\maketitle\thispagestyle{plain}
\begin{abstract}
Fix $m$ offsets and add the same radix-$B$ word to every offset.
At a cut after $i$ digits, the exact number of reachable joint carry states is
one plus the number of distinct nonzero offset residues modulo $B^i$.
The states form a chain, so their number is at most
$\min\{m+1,B^i\}$. This bound is attained and gives the minimum number of
residual states that preserve every remaining output digit.
For binary words, the ordered states support exact counting under arbitrary
input and output bit masks. Stable partition and monotone sweeps construct
all transitions in $O(wm)$ arithmetic operations. The counter also returns
the least satisfying word. Thus 32 parallel additions have at most
33 carry states at a cut, although a general vector of 32 independent carries
has $2^{32}$ values. Exact source and independent whole-input comparisons
accompany the proofs.
\end{abstract}
\noindent\textit{Keywords.} Carry propagation; shared-input addition;
residual states; finite-state counting; bit masks.

\section{Parallel sums and their residual functions}
Fix integers $B\geq2$, $w\geq1$, and $m\geq1$.
Let $\alpha_1,\ldots,\alpha_m\in\{0,\ldots,B^w-1\}$ be fixed offsets.
The variable input is one word $x\in\{0,\ldots,B^w-1\}$, shared by all sums
\begin{equation}\label{eq:sums}
 y_j(x)=(x+\alpha_j)\bmod B^w,\qquad 1\leq j\leq m.
\end{equation}
Each addition starts with carry zero. We process the input digits from least
to most significant. At a cut $i$, put
\begin{equation}\label{eq:carry}
 M_i=B^i,\qquad a_{j,i}=\alpha_j\bmod M_i,\qquad
 u=x\bmod M_i,\qquad
 c_{j,i}(u)=\left\lfloor\frac{u+a_{j,i}}{M_i}\right\rfloor.
\end{equation}
Every carry is either zero or one, at every radix.
The joint carry vector is $C_i(u)=(c_{1,i}(u),\ldots,c_{m,i}(u))$.

For $i<w$, a lower prefix $u$ defines a residual function of the upper
input word $v\in\{0,\ldots,B^{w-i}-1\}$:
\begin{equation}\label{eq:residual}
 R_{i,u}(v)_j=
 \left(\left\lfloor\frac{\alpha_j}{M_i}\right\rfloor
       +v+c_{j,i}(u)\right)\bmod B^{w-i}.
\end{equation}
This function returns the upper $w-i$ digits of every sum.
The lower output digits are already consumed. A residual state must retain
exactly the distinctions that can affect the remaining outputs.

\section{Exact carry width}
Let $d_i$ be the number of distinct nonzero members of
$\{a_{1,i},\ldots,a_{m,i}\}$. Let $q_i$ be the number of joint carry
vectors attained by $u\in\{0,\ldots,M_i-1\}$.

\begin{theorem}\label{thm:width}
For every cut $0\leq i\leq w$,
\begin{equation}\label{eq:width}
 \boxed{q_i=1+d_i\leq\min\{m+1,B^i\}.}
\end{equation}
The vectors $C_i(u)$ form a chain under componentwise order as $u$ increases.
For every $i<w$, these vectors are exactly the equivalence classes of the
residual functions in~\eqref{eq:residual}. Consequently, $q_i$ states are
necessary and sufficient to retain every future output.
\end{theorem}

\begin{proof}
An offset with $a_{j,i}=0$ has carry zero for every prefix.
Every other offset has one threshold,
\begin{equation}\label{eq:threshold}
 \theta_{j,i}=M_i-a_{j,i},\qquad
 c_{j,i}(u)=\ind\{u\geq\theta_{j,i}\}.
\end{equation}
The distinct thresholds lie in $\{1,\ldots,M_i-1\}$.
Write them as $\theta_1<\cdots<\theta_{d_i}$, and add
$\theta_0=0$, $\theta_{d_i+1}=M_i$.
The carry vector is constant on each nonempty integer interval
\[
 I_{i,r}=[\theta_r,\theta_{r+1})\cap\mathbb Z,
 \qquad 0\leq r\leq d_i.
\]
At each internal threshold, at least one component changes from zero to one.
No component changes back. The $d_i+1$ intervals therefore have different
vectors, in chain order. This proves the count and upper bounds.

Equal carry vectors give equal residual functions by~\eqref{eq:residual}.
For unequal vectors, choose a component $j$ where they differ.
At $v=0$, the corresponding residual values differ by one modulo
$B^{w-i}$. Since $i<w$, this modulus is at least two, so the values differ.
The residual functions are distinct. Any deterministic summary preserving
all these responses must distinguish all $q_i$ classes.
\end{proof}

\begin{corollary}\label{cor:sharp}
The upper bound in~\eqref{eq:width} is sharp. At a prescribed cut, choose
$\min\{m,B^i-1\}$ distinct nonzero offset residues and repeat residues as
needed. At a nonterminal cut, the minimum fixed-length residual label has
$\lceil\log_2q_i\rceil$ bits.
\end{corollary}

The label count presumes the offsets and cut are known.
At the terminal cut $i=w$, the carry vectors remain as counted in
\eqref{eq:width}. If only modular outputs are required, there are no
remaining digits and all terminal states can be accepted together.
If the outgoing carry vector is retained, its $q_w$ values remain distinct.
Input restrictions can remove reachable classes. They cannot add new ones.

For two offsets, the possible vectors are $00,10,11$ or $00,01,11$,
according to their low-residue order. Equal residues leave only equal carries.
For 32 offsets, at most 33 vectors occur. Six bits suffice for their rank;
choices with 33 distinct states require six bits.
These labels describe the carry interface, not the input word or the full
vector of sum outputs.

\section{Constructing the binary transitions}
We now take $B=2$. Retain a threshold for every offset, using
$\theta_{j,i}=M_i$ when $a_{j,i}=0$.
The state label $r$ is the interval index in the proof above.
Define the switch rank $s_{j,i}$ by
\begin{equation}\label{eq:switch}
 c_{j,i}=\ind\{r\geq s_{j,i}\}.
\end{equation}
For a zero residue, set $s_{j,i}=q_i$, so that this carry is always zero.
Repeated thresholds share a switch rank.

Let $e_j$ be bit $i$ of $\alpha_j$, and let $b$ be the next input bit.
The ordinary addition equations are
\begin{equation}\label{eq:transition}
 z_j=e_j\oplus b\oplus c_{j,i},\qquad
 c_{j,i+1}=\left\lfloor\frac{e_j+b+c_{j,i}}2\right\rfloor.
\end{equation}
Here $z_j$ is output bit $i$. Thus a state and an input bit determine the
output bits and the next state.

\begin{lemma}\label{lem:stable}
Given the current ordered thresholds, the next ordered thresholds and both
state-transition maps can be constructed in $O(m)$ arithmetic operations.
\end{lemma}
\begin{proof}
The next threshold is
\begin{equation}\label{eq:threshold-update}
 \theta_{j,i+1}=\theta_{j,i}+(1-e_j)M_i.
\end{equation}
All thresholds with $e_j=1$ lie in $[1,M_i]$.
All with $e_j=0$ lie in $[M_i+1,2M_i]$.
Within each group the old order is preserved. A stable partition, with
$e_j=1$ first, therefore constructs the next order in linear time.
One scan identifies distinct internal thresholds and assigns switch ranks.
The value $2M_i$ represents a zero residue and is not an internal threshold.

For state $r$, use its left endpoint $u=\theta_r$ as a representative.
After reading $b$, the new lower input word is $u+bM_i$.
Every prefix in the same old interval has the same successor by
\eqref{eq:transition}. The representative therefore determines its successor.
For fixed $b$, these representatives increase with $r$.
A simultaneous scan through them and the next thresholds locates every
successor in $O(q_i+q_{i+1})=O(m)$ operations.
Repeat the scan for the other input bit.
\end{proof}

At cut zero, every threshold is one and the only state is the zero vector.
Lemma~\ref{lem:stable} constructs all subsequent layers without a table of
$2^m$ carry vectors or a comparison sort at each cut.

\section{Exact counting with bit masks}
For each sum, fix a mask $\mu_j$ and a target $\tau_j$ with
$\tau_j\andb\neg\mu_j=0$. We count the words satisfying
\begin{equation}\label{eq:masks}
 y_j(x)\andb\mu_j=\tau_j\quad(1\leq j\leq m),\qquad
 x\andb\mu_0=\tau_0.
\end{equation}
The input mask is optional; zero imposes no input restriction.
All masks and targets have width $w$.

For a proposed next bit $b$, every specified output bit requires
\begin{equation}\label{eq:required}
 c_{j,i}=e_j\oplus b\oplus (\tau_j)_i.
\end{equation}
A required carry one means $r\geq s_{j,i}$, and a required carry zero means
$r<s_{j,i}$. The allowed state labels form one interval, possibly empty:
\begin{equation}\label{eq:interval}
 \begin{aligned}
 \ell_b&=\max\bigl(\{0\}\cup\{s_{j,i}:c_{j,i}\text{ is required to be }1\}\bigr),\\
 h_b&=\min\bigl(\{q_i-1\}\cup\{s_{j,i}-1:c_{j,i}\text{ is required to be }0\}\bigr).
 \end{aligned}
\end{equation}
Only masked output positions contribute to these sets.
The input mask can suppress either proposed bit altogether.

\begin{theorem}\label{thm:count}
The solution count of~\eqref{eq:masks}, and the least satisfying word when
one exists, can be computed in $O(wm)$ operations on integers of
$O(w+\log(m+1))$ bits. There are at most
\begin{equation}\label{eq:attempts}
 2\sum_{i=0}^{w-1}q_i\leq2w(m+1)
\end{equation}
state-transition attempts. Streaming storage uses $O(m)$ such integers,
including the input description. A conservative elementary bit-operation
bound is $O(wm(w+\log(m+1)))$.
\end{theorem}

\begin{proof}
Maintain one exact counter $n_i(r)$ per state.
Its value is the number of lower input prefixes that satisfy all restrictions
in positions below $i$ and reach state $r$.
Initially $n_0(0)=1$.
For each $b\in\bits$, compute the interval~\eqref{eq:interval} and the
successor map from Lemma~\ref{lem:stable}.
For every allowed $r$, add $n_i(r)$ to its successor's counter.

Each prefix has exactly one state and exactly one next input bit.
Every accepted extension satisfies the new mask bits by
\eqref{eq:required}. Conversely, every legal extension passes this test.
This proves the counter invariant by induction, with neither omission nor
duplicate counting. Summing the terminal counters gives the exact count.

Alongside each positive counter, retain the least prefix attaining it.
For a fixed state, all such prefixes have the same possible suffixes.
Replacing a prefix by the least one preserves validity and cannot increase
the resulting complete word. Each transition extends this stored prefix by
$b2^i$, and a successor retains the smaller of the candidates it receives.
Induction proves that the minimum terminal witness is the least solution.

Each layer needs a constant number of scans over offsets and states.
The threshold construction, mask bounds, and two transitions therefore cost
$O(m)$ arithmetic operations. Counters and witnesses have at most $w+1$ bits;
thresholds have at most $w+1$ bits, and indices need $O(\log(m+1))$ bits.
Two counter layers, two witness layers, and the threshold arrays suffice.
The stated time and storage bounds follow. Recording all layer traces is an
optional additional output and uses additional storage.
\end{proof}

The constructor and counter are included in the bound. It is not a bound
only for querying an already built graph. The mask constraints are local to
individual bit positions; unrestricted additional predicates need their own
state and evaluation cost. The result is a counting and witness algorithm,
not an enumeration of every satisfying word.

\section{Examples and exact verification}
At cut eight, choose offsets $(13,45,93,157,221)$ in nine-bit words.
The five thresholds are $35,99,163,211,243$.
All 256 lower prefixes lie in six classes of sizes
\[
 35,\quad64,\quad64,\quad48,\quad32,\quad13.
\]
In the given offset order, the carry vectors are
$00000,00001,00011,00111,01111,11111$.
Every vector gives a different final-digit response. The figure source
records every prefix-to-class incidence.

A separate fixed example uses 32 offsets at width 256 and at most two specified
output positions per offset. The exact solver counts
\[
 291507445114675457594661599285981859699083240676748846694400
\]
solutions. It constructs 8,290 nonterminal source states and attempts
16,580 transitions; its maximum layer width is 33.
A separately written ripple counter, whose keys are full carry vectors,
returns the same count and least witness. Complete offsets, masks, targets
and witness are supplied in \texttt{evidence/verification.json}.
The example illustrates the proved representation and counting bounds;
it supplies no comparison of processor speed or electrical cost.

The finite verification uses direct integer addition as its small-domain
oracle. It checks 30,792 prefix panels and 211,326 prefixes, including
repeated and zero offset residues. First-suffix outputs distinguish all
41,560 checked pairs of residual states. In 3,840 complete masked queries,
all 785,664 input words are evaluated directly. Counts and least witnesses
agree with both implementations. A further 168 queries cover word widths
16 through 1,024 and two through 64 parallel sums. Every comparison passes.
The general results follow from the preceding proofs; these checks qualify
the implementation rather than replace those proofs.

\section{Relation to existing methods}
Ordered decision diagrams merge prefixes with identical residual
behavior~\cite{bryant}. Finite automata for arithmetic constraints on binary
representations are also established~\cite{wolper}.
The present construction specializes those methods to fixed translations
of one shared word. Its explicit threshold order gives the exact layer
width, a matching residual-state lower bound, and a uniform linear-layer
construction for masked counting.

The common-input assumption is central. Independent input words do not
have the common threshold variable in~\eqref{eq:threshold}.
Cascaded addition with independently variable incoming carries has a
different response-function classification~\cite{cascade}.
Here the offsets are parameters, and one word is the common variable.
The stated bounds concern this arithmetic family and its specified outputs.

\section*{Reproducibility}
\texttt{frontier.py} contains the ordered solver. \texttt{verify.py} contains
the direct-input oracle, the independent full-vector counter, and the
whole-input experiments. \texttt{visual/build\_scene.py} builds the diagram
from \texttt{evidence/figure.json} in native Blender.
The manuscript and figure are licensed CC BY 4.0; the executable source
is licensed MIT. The source and figure accompany the DOI-linked release.

\begin{thebibliography}{9}
\bibitem{bryant}
R. E. Bryant, Graph-based algorithms for Boolean function manipulation,
\textit{IEEE Transactions on Computers} C-35(8) (1986), 677--691.
\href{https://doi.org/10.1109/TC.1986.1676819}{doi:10.1109/TC.1986.1676819}.
\bibitem{wolper}
P. Wolper and B. Boigelot, An automata-theoretic approach to Presburger
arithmetic constraints, in \textit{Static Analysis}, LNCS 983,
Springer, 1995, pp. 21--32.
\href{https://doi.org/10.1007/3-540-60360-3_30}{doi:10.1007/3-540-60360-3\_30}.
\bibitem{cascade}
lmtlssss{\authorsymbol 🜐}, \textit{Exact Counts of Sum-and-Carry Functions
for Cascaded Addition}, 2026.
\href{https://doi.org/10.5281/zenodo.23700193}{doi:10.5281/zenodo.23700193}.
\end{thebibliography}
\end{document}
